\documentclass[11pt]{article}
\usepackage[T1]{fontenc}
\usepackage[utf8]{inputenc}
\usepackage{lmodern}
\usepackage[margin=1in]{geometry}
\usepackage{microtype}
\usepackage[hyphens]{url}
\usepackage{hyperref}
\setlength{\parskip}{0.45em}
\setlength{\parindent}{0pt}
\setlength{\emergencystretch}{3em}
\title{Learning by Viewing -- Kickoff Pack}
\date{}
\begin{document}
\maketitle
\section*{1. Project overview}
Core question: does scalable access to high-quality video demonstration (for example, YouTube) causally increase skill acquisition relative to text, diagrams, and trial-and-error learning alone?\\

Primary research questions:\\
1. How much does video access change the probability, speed, and quality of mastering practical skills?\\
2. Which skills benefit most from viewing (procedural, motor, tacit, sequence-heavy) versus reading/diagrams?\\
3. Does video learning reduce or widen skill inequality across income, education, age, and geography?\\
4. What is the welfare impact when video substitutes for paid instruction, apprenticeships, or formal training?\\
5. Do platform dynamics (creator competition, recommendation systems, feedback loops) improve teaching quality over time?\\

Scope and framing:\\
\begin{itemize}
\item Main paper should be empirical economics first, with reduced-form and quasi-experimental evidence.
\item Theory/modeling should discipline mechanism interpretation and generate heterogeneity predictions.
\item Minor extension: interactive human/AI help as a complement to passive viewing.
\end{itemize}

Candidate outcomes:\\
\begin{itemize}
\item Individual: task completion, error rate, time-to-completion, certification/passing, wages in skill-linked occupations.
\item Household: DIY spending efficiency, contractor substitution, repair quality/safety incidents.
\item Market: service prices, demand for professional instruction, entry into skilled occupations/hobbies.
\item Distributional: gains by baseline skill and digital access.
\end{itemize}

Expected contribution:\\
\begin{itemize}
\item Introduce and measure "learning by viewing" as distinct from learning by doing and learning by reading.
\item Provide causal evidence on when demonstration-rich media improves human capital formation.
\item Quantify welfare and distributional implications in a platform-mediated learning environment.
\end{itemize}

\section*{2. Mechanisms and testable hypotheses}
Main hypotheses:\\
1. Video access raises skill acquisition rates more for procedural/motor tasks than for conceptual tasks.\\
2. Effects are larger when learners can sort across many instructors and formats (choice set mechanism).\\
3. Video narrows geographic inequality in skill access but can widen digital-access inequality.\\
4. Platform learning effects grow over time as instructional best practices diffuse among creators.\\
5. Video + minimal interactivity (comments, Q\&A, AI tutor) outperforms video-only for troubleshooting-heavy tasks.\\

Key mechanisms/channels:\\
1. Demonstration fidelity: moving visual sequences transmit tacit steps that text/diagrams miss.\\
2. Lower search/matching costs: users find instructor style/pace that fits prior knowledge.\\
3. Repetition and pacing control: pause/rewind/slow-motion increases effective practice intensity.\\
4. Quality competition/network effects: creators optimize pedagogy under feedback and engagement.\\
5. Scale and long tail supply: niche skills receive tutorials unavailable in local markets.\\
6. Social proof and motivation: visible success examples increase effort and persistence.\\
7. Complementarity with doing: viewing reduces early mistakes, making practice more productive.\\

Testable predictions:\\
1. Treatment effects should be strongest for tasks requiring ordered hand movements and error-sensitive sequencing.\\
2. Effects should increase with broadband quality/video load speed and decrease with outages.\\
3. Gains should be larger where local in-person instruction is scarce/high-cost.\\
4. As platform maturity rises, average tutorial quality metrics should improve and learning effects should strengthen.\\
5. Video-heavy learners should show fewer "first-try" errors and shorter completion times.\\

Expected welfare implications:\\
\begin{itemize}
\item Likely winners: remote/rural learners, low-income self-learners, career-switchers, small creators of instructional content, consumers via cheaper services.
\item Potential losers: incumbent local instructors with low differentiation, low-digital-access groups, workers in tasks where amateur substitution reduces paid demand.
\item Inequality possibilities:
\item Democratization channel: low-cost access broadens skill acquisition.
\item Divergence channel: those with devices, bandwidth, and digital literacy capture larger gains.
\item Net welfare sign likely positive if quality/safety externalities are managed and access gaps are addressed.
\end{itemize}

\section*{3. Modeling menu (34 options)}
1. Binary mastery choice model.\\
\begin{itemize}
\item Agents/choices/timing/state: individual chooses learning mode (view/read/do) before attempting task; state is baseline ability.
\item What viewing does: increases success probability and lowers time cost.
\item Comparative statics/predictions: larger viewing premium for low baseline ability in procedural tasks.
\item Data mapping: completion rates, attempt durations, modality exposure logs.
\end{itemize}

2. Sequential learning-with-errors model.\\
\begin{itemize}
\item Agents choose modality each period; state is accumulated mistakes.
\item Viewing reduces transition probability into high-error states.
\item Predicts fewer catastrophic failures early in learning.
\item Map with repeated attempt data and error-type classifications.
\end{itemize}

3. Bayesian signal precision model.\\
\begin{itemize}
\item Learner updates belief about correct method from noisy signals.
\item Viewing provides higher-precision signal than text.
\item Predicts larger gains when prior uncertainty is high.
\item Use pre/post confidence and objective test scores.
\end{itemize}

4. Search and matching over instructors.\\
\begin{itemize}
\item Learner samples creators with heterogeneous teaching fit.
\item Viewing expands choice set and lowers sampling cost.
\item Predicts gains from larger platform penetration and recommendation quality.
\item Map using clickstream diversity, watch histories, completion outcomes.
\end{itemize}

5. Dynamic discrete choice of effort.\\
\begin{itemize}
\item Each period learner allocates time across viewing/practice/reading.
\item Viewing raises marginal return to practice.
\item Predicts complementarity: viewing and doing increase together.
\item Use time-use + progression outcomes.
\end{itemize}

6. Human capital production with media inputs.\\
\begin{itemize}
\item Skill stock evolves with inputs: practice, text, video.
\item Video has task-specific productivity parameter.
\item Predicts modality elasticities vary by task type.
\item Estimate with panel on practice intensity and outcomes.
\end{itemize}

7. Hazard model of time-to-mastery.\\
\begin{itemize}
\item State is not-yet-mastered; exit hazard depends on modality exposure.
\item Viewing shifts hazard upward.
\item Predicts shorter median time-to-mastery with video access.
\item Map with cohort onboarding and milestone timestamps.
\end{itemize}

8. Multi-task allocation model.\\
\begin{itemize}
\item Learner chooses which skills to attempt under budget.
\item Viewing lowers fixed learning cost for some tasks.
\item Predicts broader skill portfolio and niche-skill entry.
\item Measure number/type of new skills adopted.
\end{itemize}

9. Peer diffusion model.\\
\begin{itemize}
\item Individuals observe peers' success and adopt modalities.
\item Viewing adoption diffuses through networks.
\item Predicts local peer exposure amplifies individual treatment effects.
\item Use social graph/proxy and staggered adoption.
\end{itemize}

10. Two-sided platform quality model.\\
\begin{itemize}
\item Creators choose instructional quality; learners choose watch time.
\item Better pedagogy attracts demand; feedback improves quality.
\item Predicts rising average instructional quality over platform maturity.
\item Data: video metadata, retention, outcome-linked surveys.
\end{itemize}

11. Tournament among creators.\\
\begin{itemize}
\item Creator payoff depends on ranking/engagement.
\item Viewing ecosystem induces innovation in teaching format.
\item Predicts format convergence toward high-effectiveness designs.
\item Use content feature trends and performance metrics.
\end{itemize}

12. Attention constraint model.\\
\begin{itemize}
\item Learner has finite cognitive bandwidth.
\item Viewing can reduce extraneous cognitive load for some tasks.
\item Predicts non-monotonic effects for overly edited/fast videos.
\item Map with pace/length features and outcomes.
\end{itemize}

13. Information overload/search friction model.\\
\begin{itemize}
\item Large content supply creates selection frictions.
\item Viewing benefit depends on recommendation/filter quality.
\item Predicts stronger effects where curation improves.
\item Use algorithm changes or playlist tools as shocks.
\end{itemize}

14. Mistake-cost minimization model.\\
\begin{itemize}
\item Learner chooses modality to minimize expected mistake damage.
\item Viewing most valuable when mistakes are costly.
\item Predicts strongest demand in high-penalty tasks (electrical, medical prep).
\item Map with task risk ratings and modality usage.
\end{itemize}

15. Local instruction substitution model.\\
\begin{itemize}
\item Household chooses between paid instructor and self-learning via video.
\item Viewing lowers demand for paid instruction.
\item Predicts price/demand effects in local lesson markets.
\item Use service market prices and platform penetration.
\end{itemize}

16. Complementarity with formal education model.\\
\begin{itemize}
\item Student chooses extra video supplement to classroom learning.
\item Viewing boosts returns to formal training.
\item Predicts bigger gains in under-resourced institutions.
\item Map with school-level resources and online exposure.
\end{itemize}

17. Option value model for experimentation.\\
\begin{itemize}
\item Viewing lowers uncertainty before starting skill.
\item More people attempt skills with positive option value.
\item Predicts increased trial of previously intimidating tasks.
\item Use new-channel/topic entry and novice participation metrics.
\end{itemize}

18. Quality threshold model.\\
\begin{itemize}
\item Skill mastery requires crossing an instruction-quality threshold.
\item Viewing increases chance of encountering above-threshold instruction.
\item Predicts discrete jump in mastery for users exposed to top-tail videos.
\item Map with creator quality rankings and learner outcomes.
\end{itemize}

19. Heterogeneous learning-style model.\\
\begin{itemize}
\item Learners differ in visual/verbal preference parameter.
\item Viewing effects heterogeneous by preference/cognitive traits.
\item Predicts sorting and differential gains.
\item Use baseline assessments + modality interactions.
\end{itemize}

20. Learning depreciation/refresh model.\\
\begin{itemize}
\item Skill decays without reinforcement.
\item Viewing enables low-cost refresher episodes.
\item Predicts slower skill decay and better retention.
\item Map with follow-up tests and refresh watch behavior.
\end{itemize}

21. Credential production model.\\
\begin{itemize}
\item Workers choose learning paths to pass skill exams.
\item Viewing raises pass probability at low cost.
\item Predicts exam attempt and pass-rate increases after access shocks.
\item Use certification data by region/time.
\end{itemize}

22. Spatial access model.\\
\begin{itemize}
\item Remote regions face high cost of in-person training.
\item Viewing compresses distance penalty.
\item Predicts larger treatment effects in low-density areas.
\item Map with rural broadband rollout and local training supply.
\end{itemize}

23. Household bargaining model.\\
\begin{itemize}
\item Household allocates time/resources to member upskilling.
\item Viewing lowers monetary cost and schedule coordination burden.
\item Predicts stronger gains where time constraints bind (parents/caregivers).
\item Use household composition interactions.
\end{itemize}

24. Firm training adoption model.\\
\begin{itemize}
\item Firms choose between in-house trainer and video modules.
\item Viewing lowers training cost per worker.
\item Predicts SME adoption and productivity gains in routine technical tasks.
\item Map with firm surveys and output/productivity measures.
\end{itemize}

25. Apprenticeship queue model.\\
\begin{itemize}
\item Limited apprentice slots constrain doing-based learning.
\item Viewing partially substitutes for queue-constrained slots.
\item Predicts higher entry into constrained trades.
\item Use waitlist data and training enrollments.
\end{itemize}

26. Safety externality model.\\
\begin{itemize}
\item Learners choose risky DIY with/without video preparation.
\item Viewing may reduce per-attempt risk but increase attempts.
\item Predicts ambiguous net accident levels; clearer effect on severity.
\item Map with injury rates, severity, and DIY activity proxies.
\end{itemize}

27. Quality-adjusted output model.\\
\begin{itemize}
\item Task output has quality index, not just completion.
\item Viewing increases precision and finish quality.
\item Predicts quality gains even when completion rates unchanged.
\item Use expert ratings, returns/complaints, rework rates.
\end{itemize}

28. Multi-stage production task model.\\
\begin{itemize}
\item Task has stages; failure at early stage cascades.
\item Viewing especially improves stage transitions.
\item Predicts largest gains on bottleneck stages.
\item Map with step-level telemetry or process audits.
\end{itemize}

29. Algorithmic exposure model.\\
\begin{itemize}
\item Platform recommender determines instructional exposure quality.
\item Viewing effect mediated by recommendation policy.
\item Predicts discontinuities around algorithm updates.
\item Use policy change/event study with creator-learner panels.
\end{itemize}

30. Creator entry model.\\
\begin{itemize}
\item Potential instructors decide whether to produce tutorials.
\item Lower production tools + demand induce entry.
\item Predicts more supply in underserved skills and faster pedagogic innovation.
\item Map with category-level creator entry and learner outcomes.
\end{itemize}

31. Language accessibility model.\\
\begin{itemize}
\item Multilingual/subtitled videos reduce language barriers.
\item Viewing lowers comprehension cost for migrants/non-native speakers.
\item Predicts larger gains in language-mismatched populations.
\item Use subtitle availability and migrant outcomes.
\end{itemize}

32. Social reputation model.\\
\begin{itemize}
\item Learners post outcomes publicly; reputation rewards mastery.
\item Viewing strengthens norm of self-teaching and sharing.
\item Predicts peer-driven adoption cycles and persistence.
\item Map with social posting data and subsequent learning behavior.
\end{itemize}

33. Income-constrained investment model.\\
\begin{itemize}
\item Learning investment constrained by cash and liquidity.
\item Viewing is low-cost capital-light learning input.
\item Predicts larger effects among credit-constrained groups.
\item Use income/credit proxies and heterogeneity.
\end{itemize}

34. Hybrid video + AI helper model (extension).\\
\begin{itemize}
\item Learner chooses passive video only or video plus interactive assistant.
\item Interactivity lowers troubleshooting friction after failed attempts.
\item Predicts gains concentrated in tasks with high idiosyncratic failure modes.
\item Map with usage logs of Q\&A/assistant tools and retry outcomes.
\end{itemize}

\section*{4. Literature map (what to search + why)}
Bucket A. Economics of learning / skill acquisition / human capital\\
\begin{itemize}
\item Subtopics:
\item Learning-by-doing and experience curves.
\item Observational/social learning and peer effects.
\item Skill formation and dynamic human capital accumulation.
\item Tacit knowledge transfer and apprenticeship economics.
\item Keywords:
\item "learning by doing economics"
\item "observational learning human capital"
\item "social learning labor economics"
\item "tacit knowledge transfer productivity"
\item "experience curve worker performance"
\item What we need from this bucket:
\item Theoretical foundations for mechanisms.
\item Established predictions on heterogeneity and dynamic gains.
\item Benchmarks for comparing viewing to doing/reading.
\end{itemize}

Bucket B. Economics of online platforms / information / media / education technology\\
\begin{itemize}
\item Subtopics:
\item Platform content quality competition.
\item Returns to online education and digital learning tools.
\item Attention markets and recommendation systems.
\item Video-based instruction versus text-based instruction.
\item Keywords:
\item "online platform content quality economics"
\item "video learning outcomes quasi experimental"
\item "education technology causal effects"
\item "YouTube education skill acquisition"
\item "recommendation algorithm learning outcomes"
\item What we need from this bucket:
\item Evidence on platform-mediated information quality and access.
\item Empirical designs that isolate platform exposure.
\item Measurement conventions for digital learning behavior.
\end{itemize}

Bucket C. Identification strategies used in related work\\
\begin{itemize}
\item Subtopics:
\item Broadband rollout and internet speed shocks.
\item Platform outages/blocks and access interruptions.
\item Device adoption shocks (smartphone penetration, data pricing).
\item Differential timing of local digital infrastructure upgrades.
\item Keywords:
\item "broadband rollout difference in differences skills"
\item "internet outage natural experiment education"
\item "platform ban block causal effects"
\item "fiber rollout event study labor outcomes"
\item "data price shock digital learning"
\item What we need from this bucket:
\item Credible treatment variation sources.
\item Threat-mitigation playbooks (pre-trends, migration, anticipation, compositional shifts).
\item Templates for event-study and IV diagnostics.
\end{itemize}

\section*{5. Empirical design brainstorm (10+)}
1. Broadband rollout x skill intensity DiD.\\
\begin{itemize}
\item Outcomes: certification rates, DIY service substitution, skill-related earnings.
\item Treatment/proxy: local high-speed internet availability.
\item ID: staggered DiD/event study across regions.
\item Threats/fixes: endogenous rollout; use planned infrastructure timing controls and pre-trend checks.
\item Feasibility: high (public infrastructure and labor data often available).
\end{itemize}

2. YouTube outage natural experiment.\\
\begin{itemize}
\item Outcomes: search-to-completion of tasks, course completion, repair outcomes.
\item Treatment: temporary outages or throttling episodes.
\item ID: event study around outage windows.
\item Threats/fixes: concurrent shocks; use unaffected platforms/regions as controls.
\item Feasibility: medium (needs granular outage and platform usage data).
\end{itemize}

3. Mobile data price shock IV.\\
\begin{itemize}
\item Outcomes: instructional video watch time, task success proxies.
\item Treatment: effective video affordability.
\item ID: IV using telecom tariff changes.
\item Threats/fixes: pricing endogeneity; exploit regulatory or exogenous policy shifts.
\item Feasibility: medium.
\end{itemize}

4. School/workplace video-access policy change.\\
\begin{itemize}
\item Outcomes: pass rates, productivity metrics.
\item Treatment: introduction of allowed/blocked video learning tools.
\item ID: DiD across institutions adopting at different times.
\item Threats/fixes: policy chosen where problems biggest; include institution trends and matched controls.
\item Feasibility: medium-high.
\end{itemize}

5. Recommendation algorithm change.\\
\begin{itemize}
\item Outcomes: learning completion, retention, follow-through behavior.
\item Treatment: exogenous shifts in exposure to instructional content.
\item ID: event study or diff-in-discontinuities using policy change date.
\item Threats/fixes: simultaneous UI changes; isolate affected categories and placebo categories.
\item Feasibility: medium.
\end{itemize}

6. Geographic instrument: terrain-driven broadband quality.\\
\begin{itemize}
\item Outcomes: skill uptake and quality.
\item Treatment: realized streaming quality.
\item ID: IV with topography/cable distance interactions.
\item Threats/fixes: exclusion concerns; saturate with geographic controls and historical development trends.
\item Feasibility: medium.
\end{itemize}

7. Platform entry of localized language subtitles.\\
\begin{itemize}
\item Outcomes: learning gains among non-native speakers.
\item Treatment: subtitle availability by language/time.
\item ID: DiD by language community exposure.
\item Threats/fixes: concurrent migration trends; add demographic and area-year controls.
\item Feasibility: medium-high.
\end{itemize}

8. Public library/community center Wi-Fi expansion.\\
\begin{itemize}
\item Outcomes: skill course completion, certification, employment transitions.
\item Treatment: nearby free high-speed access.
\item ID: DiD or synthetic control at municipality level.
\item Threats/fixes: targeted rollout; propensity-score weighting and pre-trend balance.
\item Feasibility: high.
\end{itemize}

9. Vocational exam reform requiring practical competencies.\\
\begin{itemize}
\item Outcomes: pass rates in practical components.
\item Treatment: differential need for demonstration learning.
\item ID: triple-difference (before/after x practical-heavy exams x high-video-access regions).
\item Threats/fixes: other reform components; isolate practical-specific sections.
\item Feasibility: medium.
\end{itemize}

10. Device shock: low-cost smartphone rollout.\\
\begin{itemize}
\item Outcomes: first-time engagement with skill learning and outcomes.
\item Treatment: smartphone/video-capable access.
\item ID: DiD with staggered retail rollout or subsidy eligibility thresholds.
\item Threats/fixes: income trends confounding; include local labor and price controls.
\item Feasibility: medium-high.
\end{itemize}

11. Category-specific creator supply shock.\\
\begin{itemize}
\item Outcomes: skill adoption in affected categories.
\item Treatment: sudden increase in tutorial supply from creator entry wave.
\item ID: Bartik-style exposure or event study.
\item Threats/fixes: demand-driven entry; instrument using creator-side shocks (policy, monetization changes).
\item Feasibility: medium.
\end{itemize}

12. Platform block/unblock across countries.\\
\begin{itemize}
\item Outcomes: skill market indicators, learning outcomes.
\item Treatment: access ban/lift timing.
\item ID: synthetic control + event study.
\item Threats/fixes: macro-political confounds; narrow windows and multiple controls.
\item Feasibility: medium-low (policy complexity).
\end{itemize}

13. Firm-level adoption of video SOPs.\\
\begin{itemize}
\item Outcomes: error rates, training duration, output quality.
\item Treatment: rollout of video-based standard operating procedures.
\item ID: staggered DiD within multi-site firms.
\item Threats/fixes: contemporaneous process upgrades; collect implementation details and phased controls.
\item Feasibility: high if partnership data obtained.
\end{itemize}

14. Weather-driven stay-at-home shocks interacting with video access.\\
\begin{itemize}
\item Outcomes: home task completion, online instructional consumption.
\item Treatment: exogenous time at home x internet quality.
\item ID: shift-share panel with fixed effects.
\item Threats/fixes: weather affects many behaviors; isolate task categories plausibly linked to viewing.
\item Feasibility: medium.
\end{itemize}

\section*{6. Deep-dive: DIY/COVID and safety/accidents angle}
Causal story:\\
\begin{itemize}
\item During COVID-era stay-at-home periods, households shifted toward DIY tasks.
\item Video tutorials reduced information frictions for attempting/performing tasks.
\item Net accidents can rise or fall: more attempts raise exposure, better instruction lowers per-attempt risk.
\item Therefore identify both extensive margin (attempts) and intensive margin (risk per attempt/severity).
\end{itemize}

Outcome measures (at least 3):\\
1. DIY project incidence/intensity: purchases of task-specific materials/tools.\\
2. Safety outcomes: emergency visits for DIY-related injuries, severity, admissions.\\
3. Quality outcomes: repeat purchases (rework), contractor call-backs after DIY attempts.\\
4. Labor market outcomes (optional): home-service demand shifts/prices.\\

YouTube learning exposure measures (at least 3):\\
1. Region-time watch intensity in DIY tutorial categories.\\
2. Search/query intensity for "how to" repair/build tasks.\\
3. Exogenous video access proxies: broadband speed, data affordability, platform uptime.\\
4. Creator supply in local language/region relevant DIY content.\\

Identification strategies (at least 2):\\
1. Triple-difference:\\
\begin{itemize}
\item Pre/post COVID periods x high/low broadband regions x DIY-intensive categories versus placebo categories.
\end{itemize}
2. Instrumental variables:\\
\begin{itemize}
\item Use exogenous network-quality variation (infrastructure timing, topology) for DIY video consumption.
\end{itemize}
3. Outage/event design:\\
\begin{itemize}
\item Compare DIY outcomes around local platform disruptions.
\end{itemize}

Key confounds and mitigation:\\
1. Income shocks and unemployment changed DIY demand.\\
\begin{itemize}
\item Mitigate with local labor controls and category-placebo checks.
\end{itemize}
2. Supply chain disruptions altered project feasibility.\\
\begin{itemize}
\item Include product availability indices and SKU-level controls.
\end{itemize}
3. Healthcare utilization changed during COVID.\\
\begin{itemize}
\item Use injury severity and alternative care channels; compare against non-DIY injury controls.
\end{itemize}
4. Time-at-home and stress changed behavior broadly.\\
\begin{itemize}
\item Include mobility/time-use controls and category-specific falsification tests.
\end{itemize}

Feasibility:\\
\begin{itemize}
\item Medium-high. Data integration is nontrivial, but credible variation and measurable outcomes exist.
\item Stronger if partnered data can measure attempts (denominator) to interpret accident rates correctly.
\end{itemize}

\section*{7. Deep-dive: Metal music technicality angle}
Causal story:\\
\begin{itemize}
\item YouTube increases access to demonstrations of advanced techniques (sweep picking, polyrhythms, double-kick patterns).
\item Musicians learn/practice higher-complexity techniques faster; stylistic best practices diffuse.
\item Over time, compositions/performances may become more technically demanding.
\item Risk: compositional trends also respond to taste, production tech, and selection into genres.
\end{itemize}

Outcome measures (at least 3):\\
1. Audio/MIDI complexity metrics: tempo-adjusted note density, rhythmic irregularity, technique markers.\\
2. Performance difficulty proxies: expert-rated difficulty, transcription complexity, playthrough error rates.\\
3. Supply outcomes: share of releases exhibiting high-technical features by cohort/region.\\
4. Labor outcomes: session-musician wages/premia for technical subgenres (if available).\\

YouTube exposure measures (at least 3):\\
1. Exposure to instructional metal guitar/drum channels by region/cohort.\\
2. Timing of broadband/video quality improvements for likely musician populations.\\
3. Tutorial supply shocks: entry/sudden growth of major technique-focused creators.\\
4. Platform recommendation shifts affecting instructional music visibility.\\

Identification strategies (at least 2):\\
1. Cohort-by-region DiD:\\
\begin{itemize}
\item Younger cohorts in high-video-access regions compared with older cohorts and low-access regions.
\end{itemize}
2. Creator supply shock event study:\\
\begin{itemize}
\item Track technicality trends before/after major tutorial-supply expansions.
\end{itemize}
3. IV with internet quality:\\
\begin{itemize}
\item Instrument instructional consumption using infrastructure quality shocks.
\end{itemize}

Confounds and mitigation:\\
1. Global genre trend cycles unrelated to YouTube.\\
\begin{itemize}
\item Include genre-time fixed effects and placebo genres with low tutorial intensity.
\end{itemize}
2. Recording software/instruments improved independently.\\
\begin{itemize}
\item Control for production technology adoption proxies.
\end{itemize}
3. Selection: high-ability musicians both consume tutorials and produce technical music.\\
\begin{itemize}
\item Use plausibly exogenous access shocks and pre-trend tests.
\end{itemize}

Feasibility assessment:\\
\begin{itemize}
\item Medium-low for clean causal inference at scale.
\item Measurement is possible, but linking exposure to causal outcomes is data-intensive and identification is fragile.
\item Better as a side project or mechanism illustration unless exceptional data partnership is available.
\end{itemize}

Better alternatives to metal technicality:\\
1. Broader instrument learning outcomes (exam passes, lesson progression, upload quality metrics).\\
2. Programming skills (course completion, coding test outcomes) with clearer labor-market links.\\
3. Home repair quality/safety where objective outcomes are easier to observe.\\

\section*{8. Alternative domains and best next steps}
35 alternative skill domains more tractable than metal technicality, each with one empirical path:\\
1. Cooking techniques: use recipe-video exposure shocks and meal-prep quality proxies from grocery basket shifts.\\
2. Baking/pastry: link tutorial exposure to local bakery-entry/home-baking sales by region-time.\\
3. Home electrical basics: relate tutorial watch intensity to permit pass rates and electrical service demand.\\
4. Plumbing repair: examine changes in plumber call-outs after localized video-access improvements.\\
5. Carpentry/furniture build: combine hardware SKU purchases with project completion reviews.\\
6. Auto maintenance: track mechanic-service substitution and parts purchases after access shocks.\\
7. Bicycle repair: use bike shop service records and local tutorial exposure.\\
8. Sewing/alterations: analyze machine/thread sales and alterations-service demand displacement.\\
9. Knitting/crochet: measure pattern complexity in marketplace uploads against video exposure.\\
10. Makeup application: link tutorial trends to cosmetology enrollments and product basket complexity.\\
11. Hairstyling/barber techniques: outcomes from cosmetology practical pass rates by broadband rollout.\\
12. Nail art: complexity in posted designs and local training enrollment shifts.\\
13. Fitness form training: injury-adjusted gym outcomes with instructional-video engagement.\\
14. Yoga technique: posture quality scores from app/video-linked cohorts.\\
15. Dance choreography: skill progression in studio assessments after tutorial integration.\\
16. Language pronunciation: spoken assessment scores with subtitle-rich video access shocks.\\
17. Public speaking: improvement in judged speech metrics with instructional video assignment.\\
18. Coding/programming: coding test performance by cohort exposed to tutorial channels.\\
19. Excel/data analysis: certification pass rates versus local platform access.\\
20. CAD/3D modeling: software certification outcomes around tutorial-supply increases.\\
21. Graphic design tools: portfolio quality ratings linked to instructional content exposure.\\
22. Video editing: freelance project quality/speed after tutorial-platform adoption.\\
23. Photography lighting: contest outcomes vs exposure to tutorial creators.\\
24. Music production/mixing: track-level mix-quality metrics after educational-channel growth.\\
25. Piano skills: exam-level progression with broadband/video quality improvements.\\
26. Guitar skills (general): grade-exam and performance outcomes by tutorial exposure.\\
27. Drumming skills (general): tempo/precision improvements in standardized assessments.\\
28. Drawing/illustration: quality ratings in marketplace portfolios vs tutorial access.\\
29. 3D printing workflows: print failure rates in maker spaces after training video rollout.\\
30. Gardening/horticulture: yield/plant survival outcomes with seasonal tutorial engagement.\\
31. Wood finishing/painting: defect rates in DIY project forums after access shifts.\\
32. First-aid procedures: simulation performance after assigned video modules.\\
33. Childcare techniques: parent training outcomes linked to instructional-viewing interventions.\\
34. Elder-care home assistance: caregiver competency assessments with video-based microtraining.\\
35. Small-business digital marketing skills: ad campaign performance after tutorial exposure.\\

Best near-term paper path (recommended):\\
1. Primary empirical setting: DIY/home repair and safety-quality outcomes.\\
2. Secondary robustness setting: coding or certification-based technical skills.\\
3. Keep metal technicality as illustrative extension or appendix.\\

Immediate next steps:\\
1. Pick one main outcome family with clear denominator (attempts) and quality/safety decomposition.\\
2. Build a data inventory table: outcomes, treatment proxies, geography-time unit, likely source, access risk.\\
3. Pre-register 2-3 feasible ID designs before deep data collection.\\
4. Draft a minimal conceptual model from options 7, 14, and 26 (hazard + mistake cost + safety externality).\\

\section*{9. Assumptions and open questions}
Assumptions used here:\\
1. Video access variation can be measured with enough granularity for credible quasi-experimental designs.\\
2. At least one domain will have both learning outcomes and attempt-denominator proxies.\\
3. Platform exposure is sufficiently distinct from generic internet use.\\

Open questions to resolve before committing to design:\\
1. Unit of analysis: individual, household, region, or firm?\\
2. Which domain has the cleanest measurable outcome and least selection bias?\\
3. What exposure measure is actually obtainable (watch logs, search trends, broadband proxies)?\\
4. Can we observe attempts vs success separately to avoid misinterpreting safety outcomes?\\
5. Is there a viable partner dataset (retail, platform, exam authority, insurer, hospital)?\\
6. How global vs country-specific should the first paper be?\\
7. Which heterogeneity dimensions are central (income, geography, age, baseline digital literacy)?\\
\end{document}